\begin{mistake}{2026-04-21}{06}

\begin{problemcontent}
已知 \(A\) 是 3 阶矩阵，\(A\) 的每行元素之和为 3，且齐次线性方程组 \(Ax=0\) 有通解 \(k_1[1, 2, -2]^{\mathrm{T}} + k_2[2, 1, 2]^{\mathrm{T}}\)，\(\alpha = [1, 1, 1]^{\mathrm{T}}\)，其中 \(k_1, k_2\) 是任意常数。

(1) 证明：对任意的一个 3 维列向量 \(\beta\)，向量 \(A\beta\) 和 \(\alpha\) 线性相关。

(2) 若 \(\beta = [3, 6, -3]^{\mathrm{T}}\)，求 \(A\beta\)。
\end{problemcontent}

\begin{solutioncontent}
(1) 令 \(\xi_1 = \alpha = [1, 1, 1]^{\mathrm{T}}\)，\(\xi_2 = [1, 2, -2]^{\mathrm{T}}\)，\(\xi_3 = [2, 1, 2]^{\mathrm{T}}\)。由题意知，\(\xi_1, \xi_2, \xi_3\) 线性无关，构成 \(\mathbb{R}^3\) 的一组基。
设任意 3 维列向量 \(\beta = x_1\xi_1 + x_2\xi_2 + x_3\xi_3\)。
根据已知条件，矩阵乘法运算结果为：
\[
A\xi_1 = 3\xi_1, \quad A\xi_2 = \mathbf{0}, \quad A\xi_3 = \mathbf{0}
\]
两边同乘矩阵 \(A\) 得到：
\[
\begin{aligned}
A\beta &= A(x_1\xi_1 + x_2\xi_2 + x_3\xi_3) \\
&= x_1A\xi_1 + x_2A\xi_2 + x_3A\xi_3 \\
&= 3x_1\xi_1 = 3x_1\alpha
\end{aligned}
\]
故 \(A\beta\) 必为 \(\alpha\) 的标量倍，两者线性相关。

(2) 当 \(\beta = [3, 6, -3]^{\mathrm{T}}\) 时，代入基底表示式可得线性方程组：
\[
x_1 \begin{bmatrix} 1 \\ 1 \\ 1 \end{bmatrix} + x_2 \begin{bmatrix} 1 \\ 2 \\ -2 \end{bmatrix} + x_3 \begin{bmatrix} 2 \\ 1 \\ 2 \end{bmatrix} = \begin{bmatrix} 3 \\ 6 \\ -3 \end{bmatrix}
\]
对增广矩阵进行初等行变换解得 \(x_1 = 3, x_2 = 2, x_3 = -1\)。
代入第一问推导出的结论 \(A\beta = 3x_1\alpha\)：
\[
A\beta = 3 \times 3 \alpha = 9 \begin{bmatrix} 1 \\ 1 \\ 1 \end{bmatrix} = \begin{bmatrix} 9 \\ 9 \\ 9 \end{bmatrix}
\]
\end{solutioncontent}

\mistakeknowledge{
\begin{itemize}
 \item \textbf{核心公式/定理}：特征值定义 \(A\xi_1=3\xi_1\)；解的定义 \(A\xi_2=\mathbf{0}\)；三维空间基底的线性表出 \(\beta = x_1\xi_1+x_2\xi_2+x_3\xi_3\)。
 \item \textbf{易错点}：未能将“每行元素之和为常数”翻译为矩阵与全 1 列向量相乘的形式；面对抽象的 \(\beta\)，未想到用已知的特征向量基底进行空间替换，试图硬求矩阵 \(A\)。
 \item \textbf{关联知识}：换基底法解题思想（将目标向量在已知输出结果的基底上进行分解，利用矩阵乘法的线性性质大幅简化运算）。
\end{itemize}}

\end{mistake}
